\chapter{Criterio de cierre arquitectónico del horizonte futuro}

\section{Cierre arquitectónico}
\begin{definition}
Se llama cierre arquitectónico del horizonte futuro al estado disciplinado en el que MADELEINE puede declarar suficientemente fijada, para la fase presente, una forma estable del horizonte ulterior del sistema y de la arquitectura futura del corpus sin convertir ese cierre en final absoluto.
\end{definition}

\begin{observation}
Su función es ordenar, en un solo bloque, los criterios de cierre, clausura y consolidación que afectan tanto al lado sistema como al lado corpus.
\end{observation}

\section{Condición mínima de cierre arquitectónico}
\begin{definition}
Existe cierre arquitectónico del horizonte futuro cuando pueden declararse simultáneamente:
\begin{enumerate}
  \item una clausura editorial provisional ya lograda;
  \item un cierre operativo suficiente de Tomo III;
  \item una continuidad disciplinada entre la fase actual y las fases posteriores;
  \item una arquitectura legible entre sistema, corpus, horizonte y proyección futura;
  \item ausencia de contradicción inmediata entre conservación y apertura posterior.
\end{enumerate}
\end{definition}

\section{Conservación}
\begin{definition}
Todo cierre arquitectónico del horizonte futuro debe conservar, salvo reapertura explícita y justificada:
\begin{enumerate}
  \item Tomo I como fundamento público;
  \item Tomo II como canon intermedio;
  \item Tomo III como capa metodológica-curatorial estabilizada;
  \item la memoria editorial de la fase presente.
\end{enumerate}
\end{definition}

\section{Cierre y apertura posterior}
\begin{observation}
El cierre arquitectónico del horizonte futuro es compatible con nuevas fases, reaperturas, reintegraciones y expansiones posteriores, siempre que entren con trazabilidad y memoria suficientes.
\end{observation}

\begin{observation}
No se trata de cerrar el sistema para siempre, sino de evitar que el porvenir quede sin forma o dependa de proliferación indistinta de cierres parciales demasiado próximos entre sí.
\end{observation}

\section{Condición mínima de no cierre arquitectónico}
\begin{definition}
No debe declararse cierre arquitectónico del horizonte futuro cuando la forma del desarrollo posterior siga dependiendo de ambigüedad fuerte entre horizonte, arquitectura, porvenir, continuidad y reserva.
\end{definition}

\section{Cláusula de alcance}
\begin{observation}
Este bloque compacta los antiguos criterios terminales, supremos, generales y ulteriores del tramo final de Tomo III. No constituye todavía una clausura absoluta del futuro del sistema ni una prohibición de expansión disciplinada posterior.
\end{observation}

